\begin{ExerciseBlock}
\ExerciseHeading{}{Verständnisfragen}{wave_propagation_checks}
Entscheiden Sie, welche Aussagen zutreffen, und begründen Sie Ihre Entscheidung.
\begin{enumerate}
\item Eine mechanische Welle transportiert Masse und Energie.\TrueFalseAnswer{falsch}
\item Eine harmonische fortschreitende Welle ist zeitlich und räumlich periodisch.\TrueFalseAnswer{wahr}
\item Aus der Wellenzahl $k$ lässt sich die Wellenlänge $\lambda$ bestimmen.\TrueFalseAnswer{wahr}
\item Jede Lösung der Wellengleichung ist eine Sinusfunktion.\TrueFalseAnswer{falsch}
\item Radiowellen benötigen wie Wasserwellen ein Ausbreitungsmedium.\TrueFalseAnswer{falsch}
\item Ebene elektromagnetische Wellen im Vakuum sind Longitudinalwellen.\TrueFalseAnswer{falsch}
\item Radiowellen haben eine höhere Frequenz als sichtbares Licht.\TrueFalseAnswer{falsch}
\item Bei einer ebenen elektromagnetischen Welle im Vakuum stehen $\vec E$ und $\vec B$ senkrecht auf $\vec k$.\TrueFalseAnswer{wahr}
\item Radiowellen und sichtbares Licht sind elektromagnetische Wellen in verschiedenen Frequenzbereichen.\TrueFalseAnswer{wahr}
\item Licht breitet sich in transparentem Glas schneller aus als im Vakuum.\TrueFalseAnswer{falsch}
\item Die Intensität einer ungedämpften ebenen Welle bleibt in Ausbreitungsrichtung konstant.\TrueFalseAnswer{wahr}
\end{enumerate}
\end{ExerciseBlock}
